\chapter{N+1 维生命周期记忆体}
\chaptervoice{记忆不只回答“保存了什么”，还必须回答“它发生在主体生命史的哪里、怎样约束此刻、又如何生成尚未发生的可能”。这里的“+1”是结构化时间索引，不是新增物理维度。}

\section*{理论来路与依赖接口}
本章承接第23章的 $h$、$E$、$\Theta$ 三相功能，以及第24章对更新速率与稳定性的区分。依据《比较智能差异.md》中“过去凝固—现在切面—未来生成”的生命周期直觉，并按照 \texttt{think.md} 的约束，我们把直觉重写成可审计的时间结构。设主体从启动时刻 $t_0$ 延续至当前时刻 $t$，其生命周期记忆体不是一个静态数组，而是
\begin{equation}
 \mathcal L_t=\bigl(E_{[t_0,t)},\,h_t,\,\Theta_t,\,\mathcal P_t^+\bigr),
 \qquad
 \mathcal P_t^+=p_{\Theta_t}(x_{t:t+H},y_{t:t+H}\mid h_t,g_t,b_t).
\end{equation}
其中 $E_{[t_0,t)}$ 只含已有证据，$h_t$ 是当前工作切片，$\Theta_t$ 是跨时段可复用的生成结构，$\mathcal P_t^+$ 是有限预测域 $H$ 内的未来分布。四者的证据地位不同：过去可追溯、现在可干预、未来只能预测。

\conceptfigure{images/fig-lifecycle.png}{生命周期记忆的结构时间：过去证据、当前切片与未来生成分布}

\section{现在作为激波切面}
“现在”不是无厚度的神秘瞬间，而是系统完成一次感知、估计、决策和写入所需的有限工作窗口 $W_t=[t-\delta_-,t+\delta_+]$。它把已经发生且可登记的事件同尚未发生、只能预测的结果分开。所谓“激波切面”只强调边界两侧的信息地位突变：越过执行时刻后，候选行动中的一个成为事实，其余分支仍只是反事实。

令观测为 $o_t$、行动为 $a_t$、现实回执为 $r_{t+1}$，当前切片按
\begin{align}
 h_t&=F_h(h_{t^-},o_t,\operatorname{Recall}(E,\Theta),g_t,b_t),\\
 a_t&\sim\pi_{\Theta_t}(\cdot\mid h_t),\\
 e_{t:t+1}&=(t,h_t,a_t,r_{t+1},\mathrm{src},\mathrm{status})
\end{align}
闭合。只有收到回执后，事件才能以相应状态写入 $E$；发出行动不等于行动成功。当前窗口宽度也依任务而变：毫秒控制与多日规划可以有不同 $\delta$，但必须共享明确的时钟、因果序和提交语义。由此，“现在”是操作性边界，而非独立于系统实现的宇宙学主张。

\section{过去作为逐渐凝固的结构块}
过去首先以带时间戳和来源的事件进入 $E$，随后可能被压缩为摘要、技能或模型参数进入 $\Theta$。“凝固”因此不是事件变成不可修改物体，而是其事实层逐渐获得更强的提交性，其解释层逐渐获得跨情境稳定性。设 $e_i$ 的事实记录为 $z_i$，解释为 $q_i^{(v)}$，则版本更新只能作用于解释：
\begin{equation}
 z_i^{(v+1)}=z_i^{(v)},\qquad
 q_i^{(v+1)}=U(q_i^{(v)},E_{>i},\Theta_v),
\end{equation}
除非审计证明原记录损坏；即使纠错，也要追加更正而非静默覆盖。

凝固程度可由证据完整性、重复验证、跨情境稳定性和依赖数量共同刻画，却不能直接等同于年龄。一个刚发生但有签名回执的交易事实可能高度确定；一个被反复讲述多年的错误解释仍可能不可靠。工程上应区分原始日志、可撤销摘要和被大量策略依赖的慢结构。依赖越多，修改的迁移成本越高，但高成本不构成其为真的证据。这一边界防止“历史悠久”被偷换成“必然正确”。

\section{未来作为概率生成空间}
未来不在记忆中以事件形式存在。系统保存的是生成未来分布所需的状态、模型、目标和约束。对于候选行动序列 $a_{t:t+H-1}$，未来空间写为
\begin{equation}
 p_{\Theta_t}(x_{t+1:t+H},y_{t+1:t+H}\mid h_t,a_{t:t+H-1},c_t),
\end{equation}
其中 $c_t$ 包含资源、权限和外部条件。该分布可以多峰，必须允许“未知”质量，且预测域越长通常不确定度越大。把最大概率路径画成单一未来，会抹去尾部风险与替代方案。

模拟轨迹若写入 $E$，必须标记为 \emph{predicted} 或 \emph{counterfactual}，并记录模型版本；它不能与 \emph{observed} 事件共享无差别检索权重。现实到来后，用适当评分规则校准生成器，例如负对数得分
\begin{equation}
 \ell_{t+1}=-\log p_{\Theta_t}(y_{t+1}\mid h_t,a_t),
\end{equation}
而不是只奖励“猜中最可能类别”。因此未来区域是决策所需的概率结构，不是预先写好的命运，也不是倒置的情景记忆。

\section{$E$ 的时间区间性质}
$E$ 不是一个无序样本袋，而是定义在生命周期区间上的带偏序轨迹。连续动力学可按时间排序，但并发工具调用、分布式回执和异步观察未必具有唯一全序。更稳妥的表示是
\begin{equation}
 E_{[t_0,t)}=(V_E,\prec,\lambda),
\end{equation}
其中 $V_E$ 为事件，$\prec$ 是由时钟、消息和因果依赖确定的偏序，$\lambda$ 保存主体、来源、动作、结果与状态。若 $e_i\prec e_j$，则 $e_j$ 的生成依赖 $e_i$；时间戳接近并不自动推出因果关系。

区间查询必须声明边界和观察时钟，例如 $E_{[u,v)}$ 采用左闭右开以避免相邻窗口重复计数。对缺失区段只能标记 gap，不能用模型推测静默补成事实。恢复时，检查点 $s_k$ 与其后的幂等事件共同重建状态：
\begin{equation}
 h_t=\operatorname{Fold}(s_k,E_{[t_k,t)}).
\end{equation}
若折叠结果依事件到达顺序而变化，就必须增加因果约束或冲突合并规则。这样，生命周期连续性才是可重放、可解释的结构，而不是“数据库里有很多旧文本”。

\section{$\Theta$ 不只是过去结构}
$\Theta$ 的确由过去证据学习，却不能被定义为过去的压缩副本。它包含状态转移、技能、表示变换、价值约束和不确定性模型，其功能是对未见情境产生可检验输出。设巩固过程为
\begin{equation}
 \Theta_{t+1}=\arg\min_{\Theta}\left[
   \sum_{e_i\in B_t}w_i\ell(e_i;\Theta)
   +\lambda D(\Theta,\Theta_t)+\gamma R_{\mathrm{safe}}(\Theta)
 \right],
\end{equation}
其中回放集 $B_t$ 来自过去，但目标包含保持旧能力与安全约束；是否形成有用结构，要由留出情境和未来现实验证决定。

因此，记住每个历史细节而不能迁移，不等于拥有良好的 $\Theta$；反之，一个简洁动力模型即使不复现全部细节，也可能更好地支持预测。必须同时防止两类错误：把统计规律当作无例外定律，以及用后来形成的 $\Theta$ 反向重写原始 $E$。前者需要适用域和置信度，后者需要事件—解释分层。$\Theta$ 是面向过去证据负责、面向未来预测开放的生成结构。

\section{$\Theta$ 对未来可能性的生成}
从 $\Theta$ 到未来不是一次性“想象”，而是模型、策略和约束的组合。对每个候选计划 $\alpha$，系统生成结果分布并计算风险敏感效用：
\begin{equation}
 J_t(\alpha)=\mathbb E_{\mathcal P_t^\alpha}[U(Y)]
 -\beta\operatorname{CVaR}_{q}[-U(Y)]-C(\alpha),
 \qquad \alpha\in\mathcal A(b_t).
\end{equation}
可行集 $\mathcal A(b_t)$ 先受权限、安全和资源边界裁剪；高期望收益不能抵消禁止约束。若模型外输入概率高，系统应优先观测、求助或执行可逆试探，而非让生成器补出虚假的确定性。

生成质量至少包含区分力、校准度与行动价值三个维度。分布覆盖真实结果但过宽，可能校准却无决策价值；分布尖锐但频繁遗漏真实结果，则属于过度自信。通过替换 $\Theta$ 的局部模块并保持 $h_t$ 不变，可检验某一结构对预测和行动的因果贡献。只有当改变生成结构稳定地改变未来分布或决策代价时，才能说该结构在生命周期记忆中承担未来生成职能。

\section{时间作为记忆的内化结构}
“N+1”中的 $N$ 表示内容或状态的坐标，$+1$ 表示系统内部显式维护的结构时间 $\sigma$。它可以包含物理时间戳、因果序、主体年龄、版本和相对阶段，但不是声称状态空间之外存在一条新的物理时空轴。生命周期对象可写成纤维族
\begin{equation}
 \mathfrak M=\{(\sigma,M_\sigma):\sigma\in\Sigma\},\qquad
 T_{\sigma\to\sigma'}:M_\sigma\to M_{\sigma'},
\end{equation}
其中转换 $T$ 描述重放、巩固、遗忘和版本迁移。没有这些转换规则，仅给每条记录加日期，还没有形成结构时间。

内部时间也不必与墙钟一一对应。休眠期间物理时间流逝而内部事件稀少；高频交互中短时间可发生大量结构更新。系统必须记录二者映射，避免把“未运行”误判为“无事发生”，也避免重启后把过期预测当作当前事实。这个定义支持跨会话连续性，却不推出主观时间体验。它只是让主体能够回答：某结构何时获得、依赖哪些事件、在什么版本下有效，以及恢复后应从哪里继续。

\section{生命周期记忆曲线}
单一“记忆强度随年龄衰减”不足以描述生命周期。对事件 $e_i$，至少要区分可检索性 $A_i(t)$、事实可信度 $Q_i(t)$、结构影响 $I_i(t)$ 与依赖负债 $D_i(t)$：
\begin{equation}
 m_i(t)=\bigl(A_i(t),Q_i(t),I_i(t),D_i(t)\bigr),\qquad t\ge t_i.
\end{equation}
$A_i$ 可因干扰下降、因提示暂时上升；$Q_i$ 可因外部复核改变；$I_i$ 可能在事件细节难以召回后仍通过 $\Theta$ 保留；$D_i$ 则衡量多少现有决策依赖该结构。由此，遗忘细节不等于影响消失，容易召回也不等于可信。

工程评估应画出分层曲线而非追求一条漂亮的指数。可以用间隔测试估计检索保持，用预测误差估计结构影响，用来源审计估计可信度，并通过删除或降权实验估计因果贡献。对安全关键承诺，目标不是让 $A_i$ 自然衰减，而是转入受保护存储并定期验证。生命周期曲线因此服务于保留、复核、压缩与退役决策，而不是把所有记忆统一交给一个遗忘率。

\section{早期经验为什么具有结构放大效应}
早期经验可能产生较大影响，并非因为“越早越真实”，而是因为它位于后续学习的上游。若更新递推为 $\Theta_{k+1}=U_k(\Theta_k,e_k)$，早期事件 $e_i$ 对终态的敏感度为
\begin{equation}
 \frac{\partial\Theta_T}{\partial e_i}
 =\left(\prod_{k=i+1}^{T-1}\frac{\partial U_k}{\partial\Theta_k}\right)
 \frac{\partial U_i}{\partial e_i}.
\end{equation}
它会通过表示基底、探索策略和数据选择影响后来看到什么、怎样解释以及哪些路径被强化。这是一种路径依赖；乘积也可能收缩，所以早期影响并非必然永久放大。

若早期样本稀少、学习率高且缺少现实校准，偶然经验可能锁定偏差。缓解方法包括保持探索、来源多样化、延迟不可逆巩固、保存反例、周期性重估以及在安全边界内做反事实训练。另一方面，刻意清除全部早期结构会破坏身份谱系和已验证技能。正确目标是让早期经验可追溯、可质疑、可被新证据有界修正，而不是崇拜“第一印象”或追求无历史初始化。

\section*{本章收束：记忆拥有方向，但未来尚未发生}
N+1 生命周期记忆体把三相记忆放入结构时间：$E$ 给出带证据地位的过去区间，$h$ 给出可干预的当前切片，$\Theta$ 依据两者生成未来概率。$+1$ 是内部时间组织与转换关系，不是额外物理维度；未来是生成分布，不是事实记录。下一章引入自我 $\Psi$ 时，只能在这一可追溯生命周期上定义主体连续性，不能用一个标签替代连续性证明。

\begin{readercheck}
若系统把模拟计划写入经历库而未标注来源，会破坏哪一条证据边界？为什么一个很早发生的事件不一定具有高结构影响？请说明物理时钟、因果偏序与主体版本三种时间信息各自解决什么问题。
\end{readercheck}
\cleardoublepage
